% Figure: the current-voltage and power-voltage curves of a photovoltaic
% module, with the maximum-power point. The I-V curve holds nearly constant
% current up to the knee near the open-circuit voltage; the P-V product peaks
% at the maximum-power point, where a tracker holds the operating voltage.
% Values illustrative for a single module: short-circuit current 10 A,
% open-circuit voltage 40 V, so the P-V peak sits near 0.8 Voc and 0.9 Isc,
% a fill factor of about 0.75.
\begin{figure}[htbp]
\centering
\begin{tikzpicture}
\begin{axis}[
  width=12.5cm, height=7.5cm,
  xlabel={module voltage $V$ (V)},
  ylabel={current $I$ (A)},
  xmin=0, xmax=42, ymin=0, ymax=11,
  xtick={0,10,20,30,40},
  ytick={0,2,4,6,8,10},
  axis y line=left,
  axis x line=bottom,
  tick label style={font=\scriptsize},
  label style={font=\small},
  legend style={font=\scriptsize, at={(0.03,0.03)}, anchor=south west, draw=none},
]
% I-V curve: nearly flat at Isc = 10 A, knee near 32 V, falling to Voc = 40 V.
% Model I(V) = Isc * (1 - exp((V - Voc)/Vt)) with Vt chosen for a sharp knee.
\addplot[engblue, very thick, domain=0:40, samples=200]
  {10*(1 - exp((x-40)/3.2))};
\addlegendentry{$I$--$V$ curve}
% maximum-power point marker on the I-V curve near V = 32 V, I = 9.0 A
\addplot[engred, only marks, mark=*, mark size=2.4pt] coordinates {(32,9.0)};
\node[anchor=south east, font=\scriptsize, color=engred] at (axis cs:31.5,9.2) {MPP};
\end{axis}
% second axis for the power curve, scaled to share the frame
\begin{axis}[
  width=12.5cm, height=7.5cm,
  xmin=0, xmax=42, ymin=0, ymax=330,
  xtick=\empty,
  axis y line=right,
  axis x line=none,
  ylabel={power $P=VI$ (W)},
  ytick={0,100,200,300},
  tick label style={font=\scriptsize},
  label style={font=\small},
  legend style={font=\scriptsize, at={(0.03,0.20)}, anchor=south west, draw=none},
]
% P-V curve: V times I(V), peaking at the MPP near 288 W
\addplot[engorange, very thick, dashed, domain=0:40, samples=200]
  {x*10*(1 - exp((x-40)/3.2))};
\addlegendentry{$P$--$V$ curve}
\addplot[engred, only marks, mark=*, mark size=2.4pt] coordinates {(32,288)};
\end{axis}
\end{tikzpicture}
\caption[Current-voltage and power-voltage curves of a PV module]{The
current-voltage and power-voltage curves of a photovoltaic module. The
current holds near the short-circuit value $I_{\text{sc}}$ across most of the
voltage range and collapses through a knee toward the open-circuit voltage
$V_{\text{oc}}$, where the current falls to zero. The power, the product
$P=VI$ (orange, right axis), rises with voltage, peaks at the maximum-power
point (MPP), and falls to zero at both ends. A maximum-power-point tracker
holds the operating voltage at the peak as irradiance and temperature move
it. The ratio of the peak power to the product $V_{\text{oc}}I_{\text{sc}}$
is the fill factor, here about $0.75$. Values illustrative for a single
module with $I_{\text{sc}}=10\,\text{A}$ and $V_{\text{oc}}=40\,\text{V}$.}
\label{fig:vol04ch14:pv-iv-curve}
\end{figure}
